\section{Separation and coefficient bounds at bounded block rank}\label{sec:bounded-rank}
The unit-coefficient oracles of \cref{sec:structured-oracles} rely on
special block shapes. This section allows arbitrary blocks of bounded cycle
rank. Separation and recovery then use finite libraries determined by each
block's signed cycle normals and solve no online linear program. The
parameter dependence is large, so the result is a structural bound, not a
runtime advantage over the compressed formulations.

For $s\ge1$ define
\begin{equation}\label{eq:rank-coefficient-bound}
 H_1=1,\qquad
 H_s=\max\{1,\lfloor(s-1)^{(s-1)/2}\rfloor\}\quad(s\ge2).
\end{equation}
The first four values are $1,1,2,5$. Put $r=\max_{B\in\Blocks}r_B$.
If there are no cyclic blocks, the flow and simplex domains and bridge
equations already describe the hull, so no rank-dependent procedure is needed.

\subsection{A fixed normal matrix for every state}\label{subsec:rank-normal-matrix}
Fix a cyclic block of rank $s$. Use original chord arcs as cycle coordinates,
with the fundamental-cycle matrix $C\in\{0,\pm1\}^{E_B\times s}$ of
\cref{lem:block-factorization}, so
$x_{E_B}=v_{E_B}+C\theta$ and the chord rows are the identity.
Let $\mathcal A$ contain the distinct nonzero rows of $C$ with both signs,
and let $M$ be the matrix with row set $\mathcal A$. It is totally
unimodular by \cref{lem:cycle-tu}, has rank $s$, contains $\pm e_i$,
and has at most $3^s-1$ rows. Repeated normals are combined by defining
\begin{equation}\label{eq:rank-base-rhs}
 \beta_a=\min\bigl(\{u_e-v_e:C_e=a\}\cup
                    \{v_e:-C_e=a\}\bigr),\qquad
 T=\{\theta:M\theta\le\beta\}.
\end{equation}
The lists in this formula are nonempty by construction. Every arc of a
cyclic block lies on a cycle and so has a nonzero row of $C$; hence these
inequalities include all arc bounds. The domain $T$ is bounded, and it is
nonempty when $\Flow$ is nonempty.

For each local state $j\in I_B$, write $\lambda_j=\lambda_{B,j}$ and set
\begin{equation}\label{eq:rank-state-rhs}
 \gamma_{a j}=\min\left(
 \{\lambda_j\beta_a\}\cup
 \{z_{ej}-y_jv_e:C_e=a,\ (e,j)\in\Obs\}\cup
 \{-z_{ej}+y_jv_e:-C_e=a,\ (e,j)\in\Obs\}\right).
\end{equation}
The observation lists are empty for the merged state. Each $\gamma_{aj}$
is a concave piecewise-affine function of the original coordinates. Write
\begin{equation}\label{eq:rank-state-domains}
 Q_j=\{\theta:M\theta\le\gamma_j\},\qquad
 \bar\theta_i=x_{e_i}-v_{e_i},
\end{equation}
where $e_i$ is chord $i$. With the original domains and bridge equations
enforced, \cref{thm:block-state-reduction} reduces hull membership to
\begin{equation}\label{eq:rank-local-and-sum}
 Q_j\ne\varnothing\quad(j\in I_B),\qquad
                  \bar\theta\in\sum_{j\in I_B}Q_j
\end{equation}
for each block. At zero weight, the signed coordinate bounds in
\eqref{eq:rank-state-rhs} force $\theta=0$ or reveal an inconsistent
observation. All subsequent statements therefore include zero weights.

\subsection{Finite local feasibility certificates}\label{subsec:rank-circuits}
Let $\mathcal K$ contain the primitive generators of the extreme rays of
\begin{equation}\label{eq:rank-feasibility-cone}
 K=\{q\in\R_+^{|\mathcal A|}:M^\top q=0\}.
\end{equation}
These are the positive circuits of the row configuration.

\begin{lemma}[Unit local certificates]\label{lem:rank-circuits}
Every $q\in\mathcal K$ is the indicator vector of a minimally dependent
row set of size at most $s+1$. Moreover,
\begin{equation}\label{eq:rank-feasibility-tests}
 Q_j\ne\varnothing\quad\Longleftrightarrow\quad
             q^\top\gamma_j\ge0\quad(q\in\mathcal K).
\end{equation}
A violated test yields an affine original-coordinate cut with unit
product coefficients by selecting active branches in \eqref{eq:rank-state-rhs}.
\end{lemma}
\begin{proof}
For an extreme nonnegative dependence, the dependence space on its support
must be one-dimensional. Otherwise a linearly independent dependence could
perturb its strictly positive support in both directions while remaining
nonnegative, contradicting extremality. No proper subset is dependent:
such a dependence would provide the same perturbation unless it were a
multiple of the original vector, which has full support. Hence the support
has at most $s+1$ rows. A cofactor formula on a full-rank set of columns
of this minimally dependent row matrix expresses its primitive dependence
by nonzero minors. Total unimodularity makes every nonzero coefficient
$\pm1$, and positivity makes them all $1$.

Farkas' lemma says $M\theta\le\gamma_j$ is feasible exactly when
$q^\top\gamma_j\ge0$ for every $q\in K$. The pointed polyhedral cone
$K$ is generated by its extreme rays, proving
\eqref{eq:rank-feasibility-tests}. Enumerating row subsets of size at
most $s+1$ and retaining the one-dimensional positive minimal dependencies
therefore constructs the entire finite library.

An observation occurs only in the two opposite row normals $C_e$ and
$-C_e$. A minimal positive dependence that uses both normals is exactly
that opposite pair, and selecting the observation in both rows cancels
it. Otherwise each observation occurs at most once. Selecting
active affine minima gives an affine upper bound on
$q^\top\gamma_j$ that equals its negative value at the candidate and
is nonnegative on the hull. Its product coefficients are thus $0,\pm1$.
\end{proof}

\subsection{Edge directions and a complete support library}\label{subsec:rank-support-rays}
For each set of $s-1$ independent rows of $M$, take one orientation of
its primitive integer null vector and deduplicate. Call the resulting
set $\mathcal D$. For $s=1$ set $\mathcal D=\{1\}$.

\begin{lemma}[Uniform edge directions]\label{lem:rank-edges}
Every direction $d\in\mathcal D$ satisfies
\begin{equation}\label{eq:rank-ternary-directions}
 d\in\{0,\pm1\}^s\setminus\{0\},\qquad
                        Md\in\{0,\pm1\}^{|\mathcal A|}.
\end{equation}
The set has at most $(3^s-1)/2$ elements and includes each coordinate
direction. It contains every edge direction of every nonempty $Q_j$,
including lower-dimensional states.
\end{lemma}
\begin{proof}
The cofactor null vector of $s-1$ independent rows has entries $0,\pm1$
by total unimodularity, and at least one entry is nonzero. It is therefore
already primitive. Its scalar product with another row is, up to sign,
the full determinant obtained by adjoining that row, again $0,\pm1$.
This proves \eqref{eq:rank-ternary-directions} and the cardinality bound.
Using all but one coordinate row gives each coordinate direction.

At a relative interior point of a one-dimensional face, the tight
defining rows have rank $s-1$. This holds even when $Q_j$ is
lower-dimensional: the common nullspace of the tight rows is exactly the
direction space of the minimal face, because all nontight inequalities
remain slack under sufficiently small motions in that nullspace. Rows
tight throughout $Q_j$ are among the tight rows. Choosing $s-1$
independent tight rows gives the asserted edge direction. Point states
have no edges to check.
\end{proof}

Form the central arrangement in objective space with hyperplanes
$d^\top h=0$, $d\in\mathcal D$. Let $\mathcal R$ contain both
primitive orientations of the null vector of every $s-1$ independent
directions from $\mathcal D$; for $s=1$ use $\mathcal R=\{-1,1\}$.
Redundant members of $\mathcal R$ are harmless. The coordinate
hyperplanes occur, so every full-dimensional closed chamber is pointed.
Every extreme ray of such a chamber is generated by a member of
$\mathcal R$, since its tight hyperplanes have rank $s-1$.
Cofactors and Hadamard's determinant bound give
\begin{equation}\label{eq:rank-ray-bound}
 |h_i|\le H_s\qquad(h\in\mathcal R).
\end{equation}
Indeed each cofactor has order $s-1$ and ternary entries; dividing by
the greatest common divisor to make it primitive cannot increase it.

\begin{lemma}[Finite support characterization]\label{lem:rank-support-library}
For any finite family of nonempty bounded polytopes whose edge directions
belong to $\mathcal D$, membership in their sum is equivalent to
\begin{equation}\label{eq:rank-support-tests}
 h^\top\bar\theta\le\sum_j\sigma_j(h)\quad(h\in\mathcal R),\qquad
 \sigma_j(h)=\max_{\theta\in Q_j}h^\top\theta.
\end{equation}
This includes lower-dimensional summands and sums. The empty sum has
support zero on every ray and is the singleton zero.
\end{lemma}
\begin{proof}
In the interior of a chamber, an objective is perpendicular to no
possible edge. Its maximizing face in each summand is therefore a single
vertex: a positive-dimensional face of a polytope has an edge. Along
a line segment inside the chamber, the maximizing vertex cannot change
without a nonunique maximizer at an intermediate objective. Thus each
support function is linear on the chamber, and by continuity on its
closure. Its sum is linear there as well. Every vector in a pointed
polyhedral chamber is a nonnegative combination of its extreme-ray
generators. The inequalities on $\mathcal R$ therefore imply the support
inequality for every objective in every chamber closure, which covers
all objective space. The support characterization of a nonempty compact
convex set completes the proof. For the empty sum the same argument
gives all inequalities $h^\top\bar\theta\le0$, forcing zero.
\end{proof}

This is the classical edge-direction/zonotope refinement argument
\citep{GritzmannSturmfels1993,OnnRothblum2004,MelamedOnn2014}.
Network-flow edge directions are also studied by
\citet{OnnRothblumTangir2005}. We include the proof because we need the
statement uniformly over all state right-hand sides, including
lower-dimensional slices.

\subsection{Exact support formulas and bounded dual multipliers}\label{subsec:rank-duals}
For $h\in\mathcal R$, let $\mathcal B(h)$ contain every nonsingular
$s$-row submatrix $B$ of $M$ for which
\begin{equation}\label{eq:rank-dual-vector}
 q_B=B^{-\top}h\ge0.
\end{equation}
Rows and components have the same chosen order; extending $q_B$ by zero
gives a multiplier on all rows of $M$.

\begin{lemma}[Support bases and multiplier bound]\label{lem:rank-duals}
For every nonempty state domain,
\begin{equation}\label{eq:rank-support-formula}
 \sigma_j(h)=\min_{B\in\mathcal B(h)}q_B^\top\gamma_{B,j}.
\end{equation}
The list $\mathcal B(h)$ is nonempty, and every retained multiplier is
an integer satisfying
\begin{equation}\label{eq:rank-multiplier-bound}
                     0\le(q_B)_i\le H_s.
\end{equation}
\end{lemma}
\begin{proof}
The support LP has dual
$\min\{q^\top\gamma_j:M^\top q=h,\ q\ge0\}$. Signed coordinate
rows make it feasible for every $h$. The primal is nonempty and bounded,
so duality gives a finite attained optimum. There is an optimal solution
with linearly independent supporting columns of $M^\top$. Indeed, a
dependence on a positive support permits both signs of a small feasible
perturbation; at an optimum its objective change must be zero. Moving
until a component vanishes reduces the support, and repetition yields
independence. Extend those rows to a basis of $\R^s$; the added
multipliers are zero. This proves \eqref{eq:rank-support-formula}, also
for a degenerate or lower-dimensional primal.

Total unimodularity gives $\det B=\pm1$, so $q_B$ is integral. The
stronger magnitude bound uses the defining directions of the ray rather
than a triangle inequality on $B^{-\top}h$. Choose independent
$d_1,\ldots,d_{s-1}\in\mathcal D$ whose primitive orthogonal vector
is $h$. By \eqref{eq:rank-ternary-directions}, every $Bd_i$ is ternary,
and these vectors remain independent. The integer vector $B^{-\top}h$
is orthogonal to them and primitive, because a unimodular transformation
preserves primitive integer vectors. Their cofactor vector, divided by
its common divisor, therefore equals $B^{-\top}h$ up to sign. Hadamard's
bound proves \eqref{eq:rank-multiplier-bound}. Rank one is immediate.
\end{proof}

\begin{theorem}[Bounded-rank separation and coefficients]\label{thm:bounded-rank}
For maximum cyclic block rank $r$, the sparse hull has a finite linear
description whose flow and observed-product coefficients are integers
of magnitude at most $H_r$. The sharper bound $H_{r_B}$ applies to
inequalities from block $B$. Simplex coefficients may contain rational
network data; the bound is in the specified rational row scaling.

After observation grouping, exact membership and original-coordinate
linear separation require
\begin{equation}\label{eq:bounded-rank-complexity}
 2^{O(r^2)}\bigl(|V|+|E|+m+|\Obs|\bigr)
\end{equation}
rational arithmetic operations, including construction of finite
block libraries. No online LP is required. Ungrouped observations may
add $O(|\Obs|\log|\Obs|)$ sorting comparisons.
\end{theorem}
\begin{proof}
Use the original domains and bridge equations, the local circuit tests
\eqref{eq:rank-feasibility-tests}, and the support tests
\eqref{eq:rank-support-tests}, evaluating each support by
\eqref{eq:rank-support-formula}. Their equivalence to
\eqref{eq:rank-local-and-sum} proves completeness. Each support expression
is a minimum of affine functions, because its multipliers are nonnegative
and its endpoints are affine minima. If a support test fails, choose
a minimizing basis and active endpoint branches separately for each
state. The resulting affine majorant agrees with the support sum at
the candidate, so it is a valid violated cut.

A nonsingular row basis cannot contain opposite normals. An observation
therefore appears at most once in a state's selected support expression,
with coefficient magnitude at most $H_s$ by
\cref{lem:rank-duals}. Product indices differ between states. The
aggregate uses distinct original chord coordinates with coefficients
$h_i$, bounded by \eqref{eq:rank-ray-bound}. Local circuit cuts have
unit product coefficients, and the original domains and bridge equations
have unit flow/product coefficients. Expanding all branches gives the
asserted finite description in the row scaling of the theorem.

There are at most $3^s-1$ signed normals and $(3^s-1)/2$ edge directions.
Enumerating their subsets of size at most $s+1$, all candidate support
rays, and all normal bases takes $2^{O(s^2)}$ operations and storage.
The product of the ray and basis counts has the same bound. All library
integers and their encoding lengths depend only on $s$.

Fundamental cycles can be constructed in $O(r|E|)$ graph/arithmetic
operations. Repeated normals are combined in parameter-dependent work
per arc, and each observation updates two endpoint entries. A block
with $a_B$ observed labels then needs
$2^{O(r_B^2)}(a_B+1)$ arithmetic operations to initialize and check its
local states and evaluate all supports. Since
$\sum_Ba_B\le|\Obs|$ and there are at most $|E|$ blocks, summing
these bounds proves \eqref{eq:bounded-rank-complexity}.
\end{proof}

The online operations are rational affine evaluations, comparisons, and
sums with parameter-bounded integer multipliers. Their encoding lengths
are polynomial in the input/candidate encoding length and the library
encoding length, so the arithmetic bound does not treat integers of
unbounded length as unit-size objects. Nor does the bound imply that a
universal library is computationally attractive at moderate rank: the
parallel-path oracle can be preferable even when its block ranks are
unbounded.

\subsection{Constructive recovery without an online LP}\label{subsec:rank-recovery}
The finite supports also give a direct constructive decomposition. The
method has a larger parameter factor than separation.

\begin{theorem}[Finite-basis compact recovery]\label{thm:rank-recovery}
A rational hull point with maximum cyclic block rank $r$ admits a
compact rational decomposition into at most $m+1$ graph points, computable
without an online LP in
\begin{equation}\label{eq:rank-recovery-complexity}
 2^{O(r^3)}\bigl(|V|+|E|+m+|\Obs|\bigr)
\end{equation}
arithmetic operations, including parameter-dependent preprocessing and
with the same observation-grouping proviso as
\cref{thm:bounded-rank}. The output and intermediate rational numbers
have polynomial encoding length in the input/candidate and library
encoding lengths. Dense state-flow output has the separate cost stated
in \cref{prop:theta-recovery}.
\end{theorem}
\begin{proof}
Fix a block of rank $s$. Compute every state support on $\mathcal R$
as above, order the local states, and precompute suffix sums
$S_i(h)=\sum_{j>i}\sigma_j(h)$. At step $i$, let $t$ be the remaining
aggregate, which by induction belongs to $\sum_{j\ge i}Q_j$.
The permissible next state vector is in the nonempty polytope
\begin{equation}\label{eq:rank-recovery-polytope}
 W_i(t)=\{\theta:M\theta\le\gamma_i,\quad
                  -h^\top\theta\le S_i(h)-h^\top t
                                     \quad(h\in\mathcal R)\}.
\end{equation}
By \cref{lem:rank-support-library}, this is exactly
$Q_i\cap(t-\sum_{j>i}Q_j)$. It is bounded because it is contained
in $Q_i$. Even if lower-dimensional, a nonempty bounded polytope has
a vertex, and its active defining normals at a vertex have full ambient
rank: otherwise a nonzero null direction would allow small motions in
both directions without violating any inequality.

The normal list $N$ of \eqref{eq:rank-recovery-polytope} consists only
of the rows of $M$ and $-\mathcal R$ and is fixed before the candidate
is given. Enumerate all nonsingular $s$-row submatrices of $N$ and
precompute their inverses. For each basis, solve its equations using
the current right-hand side and test the resulting vector against every
row of \eqref{eq:rank-recovery-polytope}. At least one candidate is a
feasible vertex by the preceding paragraph. Choose the first, subtract
it from $t$, and continue. The last suffix is the singleton zero, so
the construction ends with zero remaining aggregate. This proves exact
state recovery without solving an online LP.

The normal list has $2^{O(s^2)}$ rows. Its $s$-row basis enumeration,
including all feasibility checks per step, costs $2^{O(s^3)}$ arithmetic
operations and storage. The earlier support computation is cheaper.
There are $a_B+1$ local states per block, so summing the same sparse
state-count bound as in \cref{thm:bounded-rank} proves
\eqref{eq:rank-recovery-complexity}.

Sequential division does not make encoding lengths grow exponentially in
the number of states. The fixed normal entries are
integers bounded by $H_s$. Every inverse basis is its integer adjugate
divided by a nonzero integer determinant of parameter-bounded encoding
length. All initial state endpoints and supports have a common rational
denominator whose encoding length is polynomial in the input/candidate
length and the library length. At each chosen step, a common denominator
for the remaining aggregate need only be multiplied by that step's
basis determinant. After $a_B+1$ steps its encoding length has therefore
grown additively by at most $a_B+1$ parameter-bounded terms. Numerators
also have polynomial length: state points are within their finite base
bounds, remaining aggregates are sums of such points, and rejected
basis candidates are obtained from these right-hand sides by fixed
parameter-bounded matrices. This proves the claimed intermediate and
output encoding bound.

Finally, normalize positive-weight explicit state coordinates and use
the normalized merged state as the default for its constituent global
labels. The zero-weight and unused-default conventions of
\cref{prop:theta-recovery} apply unchanged. Expand through $C$ and
use \cref{thm:block-state-reduction} to obtain graph points with the
global simplex weights. A block coordinate has $s$ entries, so defaults
and observed-label exceptions give a compact representation within the
stated bound.
\end{proof}

\subsection{Coordinate sections and necessary coefficients}\label{subsec:section-facets}
The following lemma turns a local coordinate section into a coefficient
ratio that every description must contain. It is also used in
\cref{sec:universality,sec:sp-growth}. A \emph{coordinate plane}
fixes all original coordinates except two designated coordinates $U,V$;
it does not rescale or project additional free coordinates.

\begin{lemma}[Local section forces an ambient coefficient ratio]\label{lem:section-facet}
Let $P$ be a polytope and $L$ a coordinate plane with free coordinates
$U,V$. Suppose that in an open neighborhood of a point $(U_0,V_0)$,
the section $P\cap L$ is exactly
\begin{equation}\label{eq:local-halfplane}
                 \alpha U+\beta V\le\gamma,
 \qquad \alpha U_0+\beta V_0=\gamma,\qquad
                 (\alpha,\beta)\ne(0,0).
\end{equation}
Every finite linear inequality description of $P$, together with valid
affine equations, contains an inequality whose coefficients on $U,V$
are a positive multiple of $(\alpha,\beta)$. In particular, a relative
facet description contains such an inequality. Adding any equation
valid on $P$ to a facet representative cannot change those two coefficients.
\end{lemma}
\begin{proof}
The section has two-dimensional interior near $(U_0,V_0)$. Consequently
every valid affine equation restricts to an identity on $L$: its
coefficients on $U,V$ are zero, and its fixed-coordinate constant is
zero. Restrict a finite inequality description to $L$. At the boundary
point, at least one nonconstant restricted inequality must be active.
Otherwise all nonconstant inequalities are strictly satisfied in a
neighborhood, and constant inequalities cannot produce the asserted
boundary. An active valid inequality must vanish on the local boundary
line, because feasible points occur in both tangent directions. Its
normal is therefore proportional to $(\alpha,\beta)$, and feasibility
on the interior side makes the multiple positive. The argument applies
to a finite relative facet description as well. The previously proved
zero coefficients of valid equations establish representative invariance.
\end{proof}

\subsection{Sharp rank-three examples with sparse observations}\label{subsec:sharp-k4}
Use vertices $0,1,2,3$ with every edge of $K_4$ directed toward the larger
vertex, order the arcs as $(12,13,23,01,02,03)$, and set every capacity
to one. The tree is $01,02,03$, with chord matrix
\begin{equation}\label{eq:sharp-k4-matrix}
 C=\begin{pmatrix}
 1&0&0\\0&1&0\\0&0&1\\1&1&0\\-1&0&1\\0&-1&-1
 \end{pmatrix}.
\end{equation}
The following construction uses only two explicit simplex variables and
five observed products. Unlike the source--sink example in
\cref{ex:forest-k4}, it has multiple supplies and demands.

Take reference flow and balances
\begin{equation}\label{eq:sharp-k4-balances}
 v=(1/2,1/2,1/2,1/4,1/2,1/2),\qquad
 b=Av=(-5/4,-3/4,1/2,3/2)^\top.
\end{equation}
Fix $y_1=y_2=1/3$, so the residual weight is also $1/3$, and fix
\begin{equation}\label{eq:sharp-k4-aggregate}
 \bar\theta=(1/6,1/24,1/8),\qquad
 \bar x=(2/3,13/24,5/8,11/24,11/24,1/3).
\end{equation}
Observe exactly
\begin{equation}\label{eq:sharp-k4-observations}
 \Obs=\{(12,1),(13,1),(02,1),(23,2),(01,2)\}.
\end{equation}
Leave $U=z_{12,1}$ and $V=z_{13,1}$ free, and fix the other observations
to $z_{02,1}=z_{23,2}=1/6$ and $z_{01,2}=1/12$. Put
$p=U-1/6$, $q=V-1/6$. The observations give
\begin{equation}\label{eq:sharp-k4-states}
 \theta^1=(p,q,p),\qquad \theta^2=s(1,-1,0).
\end{equation}
All state deviations on arcs other than $01$ lie in $[-1/6,1/6]$;
the interval on $01$ is $[-1/12,1/4]$. In particular, the residual
state obeys $\theta^0_1\le1/6$ and
$\theta^0_2+\theta^0_3\le1/6$. For $h=(1,1,1)$, these bounds give
$h^\top\theta^0\le1/3$, while $h^\top\theta^2=0$ and
$h^\top\bar\theta=1/3$. Every feasible point in the section therefore
satisfies
\begin{equation}\label{eq:sharp-k4-cut}
                    2p+q\ge0,\qquad 2U+V\ge1/2.
\end{equation}

Conversely, suppose $|p|,|q|<1/96$ and $w=2p+q\ge0$. Choose
\begin{equation}\label{eq:sharp-k4-witness}
 s=q/2,\qquad
 \theta^0=(1/6-p-q/2,\ 1/24-q/2,\ 1/8-p).
\end{equation}
These vectors sum to \eqref{eq:sharp-k4-aggregate}. The first state's
arc deviations are $(p,q,p,p+q,0,-p-q)$, all within its scaled bounds.
The second state's arc deviations have magnitudes at most $|q|/2$
and are zero on $01$. The residual arc deviations are
\begin{equation}\label{eq:sharp-k4-residual-arcs}
 C\theta^0=(1/6-w/2,\ 1/24-q/2,\ 1/8-p,\
              5/24-p-q,\ -1/24+q/2,\ -1/6+w/2).
\end{equation}
Here $0\le w<1/32$. The first and last entries lie in their closed
intervals $[-1/6,1/6]$. The second, third, and fifth entries lie strictly
inside those intervals. The fourth lies in $(3/16,11/48)$, contained
in $[-1/12,1/4]$. Thus all scaled bounds hold, as do the balances and
observations. Disaggregation proves that the section is exactly the
half-plane \eqref{eq:sharp-k4-cut} in the asserted neighborhood.

\begin{corollary}[Rank-three sharpness with two simplex variables]\label{cor:rank-three-sharp}
The bound $H_3=2$ in \cref{thm:bounded-rank} cannot be replaced by one,
even for a unit-capacity acyclic K4 with two explicit simplex variables
and five observed products. Every finite original-coordinate description
of this hull contains an inequality whose coefficients on $U,V$ have
ratio two.
\end{corollary}
\begin{proof}
Apply \cref{lem:section-facet} with the local normal $(-2,-1)$ at
$(U,V)=(1/6,1/6)$. An integer flow/product coefficient description
with all magnitudes at most one cannot contain this ratio, whereas
\cref{thm:bounded-rank} supplies a description with bound two.
\end{proof}

\begin{remark}[A symmetric-reference variant]\label{rem:seven-product-k4}
There is also a seven-observation construction with the same graph and
$C$, the uniform reference $v=(1/2,\ldots,1/2)$, and balances
$(-3/2,-1/2,1/2,3/2)^\top$. Use three explicit variables
$y_1=y_2=y_3=1/4$, aggregate cycle vector $(1/8,0,1/8)$, and therefore
aggregate flows $(5/8,1/2,5/8,5/8,1/2,3/8)$. Observe
\[
 (12,1),(13,1),(02,1),(12,2),(03,2),(23,3),(01,3),
\]
fix all but $U=z_{12,1},V=z_{13,1}$ to $1/8$, and put
$p=U-1/8$, $q=V-1/8$. The explicit states are
$(p,q,p)$, $t(0,1,-1)$, and $s(1,-1,0)$; every state has
$|C\theta^j|\le1/8$. The residual support in direction $(1,1,1)$
is at most $1/4$, so $2p+q\ge0$, or $2U+V\ge3/8$, is necessary.
For $|p|,|q|<1/32$ and $w=2p+q\ge0$, the choices
\[
 t=-q/2,\qquad s=q/2,\qquad
 \theta^0=(1/8-w/2,0,1/8-w/2)
\]
prove sufficiency locally. The first state's arc deviations are
$(p,q,p,p+q,0,-p-q)$, the two line-state magnitudes are less than
$1/64$, and the residual arc deviations are $(a,0,a,a,0,-a)$ with
$5/64<a=1/8-w/2\le1/8$. This variant has completely symmetric state
bounds. The five-product construction instead changes $v_{01}$ and hence
the balances $b=Av$ and the flow domain itself. This makes the residual
support face an edge instead of a point and saves one explicit line state.
Changing only the reference flow within a fixed flow domain would merely
translate its cycle-coordinate representation.
\end{remark}

\Cref{cor:rank-three-sharp} concerns the retained product coordinates
themselves and is invariant under adding the hull's affine equations. It
is therefore stronger than exhibiting a dual basis with multiplier two.
\citet[Example~2]{KhademniaDavarnia2025} exhibit a projection-cone
aggregation weight two on a different graph, with equality balances
represented by pairs of opposite inequalities. That example is an
antecedent for weighted cancellation, but it does not by itself give the
facet-ratio conclusion for equality networks.
